\section{Noisy observations and conditional screening}\label{m:sec:conditioning}

We compare the conditional electron density with the nonlinear reaction
law of Section~\ref{m:sec:barriers}. The observation noise makes a rare
data event inexpensive to impose on the eigenstate. A local quantum
comparison in the resulting state will then rule out each failure of
the desired inverse relation. We use the master kernels and deterministic
barriers of Section~\ref{m:sec:barriers}, with the small parameter now
chosen from the proposed excess charge.

Suppose, towards a contradiction, that the asserted uniform excess-charge
bound fails.  By Lemma~\ref{lem:sector-global-minimum}, choose systems at
their largest strictly bound sectors such that
\[
 K_{\rm exc}:=N-Z>0,\qquad K_{\rm exc}/M\longrightarrow\infty.
\]
Fix a universal constant $B_0>2$, to be enlarged in the final argument,
and define
\begin{equation}\label{m:eq:barrier-scales}
 r_*=(B_0M/K_{\rm exc})^{1/3}\longrightarrow0,\qquad
 r_0=2^{-n}r_*,\qquad r_j=2^jr_0\quad(0\le j\le n),
\end{equation}
where the finite integer $n\ge1$ is chosen so that $r_0<Z^{-1/3}$.
The crude bound $N\le3Z$ and $M\ge1$ imply $r_*>Z^{-1/3}$.
All subsequent smallness assertions concern this sequence of systems;
their thresholds will be universal and independent of $n$.


\subsection{The common observation construction in molecular geometry}

Set $\ell_j=r_j^{1.01}$ for $j<n$. Use the unordered arrays and likelihood
versions of Section~\ref{sec:observations}, with $N$ particles, indices
$j,\ldots,n-1$, and $\mathcal F_n$ trivial. The master kernels are those
of \eqref{m:eq:master-kernel}. Define
\[
 \mu_j(y)=\E\left[\sum_{i=1}^N\mathcal K_{x_i}(y)\mid\mathcal F_j\right],
 \qquad H_j=V_s-K*\mu_j.
\]
Lemma~\ref{lem:obs-data} gives the joint versions, fixed-system regularity,
Poisson equation, and the tower property. Here the ground state attains
$E$, so Lemma~\ref{lem:obs-tilt} gives offset
$C\ell_j^{-2}\log^5(e/P_A)$ for every event of probability $P_A>0$.

\subsection{Upper tails and a spatial cap}

The observations have two uses. First, their conditional electron density
is a continuous field to which elliptic comparisons apply. Second, an
event of these observations can be imposed at the small energy cost in
Lemma~\ref{lem:obs-tilt}. We first use that cost to exclude large positive
fields and excessive local observed counts.

For this subsection fix a numerical constant $L_1\ge1$. Its value will be
chosen after the universal field caps have been obtained. At scale
$r=r_j$, put $R_+=2r$. Constants in estimates of small probabilities may
depend on $L_1$; the field thresholds below do not.

\begin{lemma}[Positive conditional-field tail]\label{m:lem:cond-tail}
If $r/2\le a_y\le5L_1r$, then, for a universal $C$ and every $s>C$,
\begin{equation}\label{m:eq:cond-tail}
 \Prob\{H_j(y)>s a_y^{-4}\}
 \le e\exp\{-c_{L_1}(s-C)^{2/5}a_y^{-1+.004}\}.
\end{equation}
All estimates hold for sufficiently small $r_*$, uniformly in $j<n$.
\end{lemma}

\begin{proof}
Write $a=a_y$ and tilt the event in question, if it has positive
probability $P_A$. By the tower property, its average of $K*\mu_j$ is
the tilted raw average of the potential of the master smears.
Nuclei within a fixed multiple of $a$ contribute at most $Ca^{-4}$
to $V_s$. More distant electron smears have exactly their bare
potential at $y$: a contributing smear near $y$ has center within
$2t_y\ll a$. Dropping the remaining negative electron terms therefore
gives
\[
 \E[H_j(y)\mid A]\le Ca^{-4}+\E_{\psi_A}J_y.
\]
Lemmas~\ref{m:lem:screen-positive-field} and~\ref{m:lem:screen-count}, with
no initial cut, bound the last expression by
$Ca^{-4}+C\sqrt{\delta/a}$. Indeed, for small $a$, their scale quantity
is bounded by $C(a^{-7}+\delta)$. Lemma~\ref{lem:obs-tilt} and
$\ell_j\ge c_{L_1}a^{1+.01}$ imply
\[
 s-C\le C a^{7/2}\ell_j^{-1}\log^{5/2}(e/P_A)
       \le C_{L_1}a^{2.49}\log^{5/2}(e/P_A).
\]
Solving for $P_A$ proves the assertion. The deterministic threshold
comes solely from the universal local field estimates.
\end{proof}

Call $\{r\le a\le2L_1R_+\}$ the participating band. Use a grid of closed
cubes $Q$ of side $r/100$ meeting this band. Let $Q'$ be concentric
enlargements, chosen
with bounded overlap to contain every ball of radius $r$ centered in
$Q$, with an additional margin. The Lipschitz bound for $a$ ensures
that $a_Q:=\inf_{Q'}a\ge r/2$ and that all points in $Q'$ fall in the
range of Lemma~\ref{m:lem:cond-tail}. Choose smooth $0\le\chi_Q\le1$
equal to one on $Q$, supported in $Q'$, with
$|\Delta\chi_Q|\le Cr^{-2}$.

\begin{lemma}[Spatial cap and its exceptional cost]\label{m:lem:cond-cap}
There are universal thresholds $C'_{\rm cap},C_{\rm cap}$ with the
following properties. Define
\[
 \begin{split}
 Q\text{ is bad}&\quad\Longleftrightarrow\quad
       \sup_QH_j>C'_{\rm cap}a_Q^{-4},\\
 B_Q&=\1_{\{Q\text{ bad}\}}\sup_Q(H_j)_+,
 \qquad \Pi=\sum_QB_Q\chi_Q,\qquad
 \mathcal E=\bigcup_{Q\text{ bad}}Q'.
 \end{split}
\]
Then $H_j-\Pi\le C_{\rm cap}a^{-4}$ on the participating band.
For each fixed $m\ge1$ and $A'>0$,
\begin{equation}\label{m:eq:cond-cap-moments}
 \Prob(Q\text{ bad})+a_Q^{4m}\E B_Q^m
       \le C_{L_1,m,A'}r^{A'}.
\end{equation}
At each deterministic point, $\Prob(y\in\mathcal E)$ has the same
arbitrarily high polynomial bound. The cube families are finite, and
\begin{equation}\label{m:eq:cond-band-volume}
 \sum_Q|Q'|\le C_{L_1}Mr^3.
\end{equation}
All these random objects are jointly measurable in the data and the
spatial variable. For the fixed system and scale, their sizes and
supports have deterministic bounds.
\end{lemma}

\begin{proof}
Since $\Delta H_j\ge-Ca_Q^{-6}$ on $Q'$, adding a quadratic and
applying the ball submean inequality gives, for any fixed large $C'$
and a universal $C''$,
\[
 \sup_QH_j\le C''a_Q^{-4}
       +Cr^{-3}\int_{Q'}(H_j-C'a_Q^{-4})_+.
\]
The same inequality follows for distributional Laplacians by
mollification. Choose $C'$ above the threshold in
Lemma~\ref{m:lem:cond-tail}. Because $a_z/a_Q$ is bounded and
$a_z\le5L_1r$ on $Q'$, integration of that tail yields
\[
 \E(H_j(z)-C'a_Q^{-4})_+^m
       \le C a_Q^{-4m}r^{A'}
\]
for arbitrary fixed $m,A'$. Jensen's inequality for the spatial
average gives the same bound for the integral term. Taking
$C'_{\rm cap}>C''+1$, Markov's inequality gives the bad-cube
probability. On a bad cube its supremum is bounded by the integral
term plus $C''a_Q^{-4}$; the probability just proved controls the
latter contribution to every moment. This proves
\eqref{m:eq:cond-cap-moments}.

If a point lies in a bad cube, the corresponding term of $\Pi$
dominates $(H_j)_+$ there. Otherwise the good-cube bound gives the
asserted cap, since $a_Q$ and the local scale are comparable by a
universal factor. At a fixed point only a bounded number of enlarged
cubes occur, proving the estimate for $\mathcal E$.
Every participating or enlarged cube lies within distance
$C_{L_1}r$ of some nucleus, by $d_y\le La_y$.
Counting grid cubes in this union proves finiteness and
\eqref{m:eq:cond-band-volume}.
Cube suprema are measurable through a countable dense subset,
because $H_j$ is continuous. The bounds from
Lemma~\ref{lem:obs-data} also give the last assertion.
\end{proof}

\subsection{The inverse comparison and the order of constants}

We now fix the constants needed for propagation. Choose $C_{\rm inv}$
larger than the barrier upper bound in $a^{-4}$ units and than
$C_{\rm cap}$. Choose $C_{\rm top}>C_{\rm inv}$ sufficiently large and
put $h_{\rm top}(y)=C_{\rm top}a_y^{-4}$; this also exceeds $h_b$ with
a fixed margin. If $\gamma$ is the constant in
Lemma~\ref{m:lem:barrier-gap}, choose
$0<\lambda_2<\lambda_1<\gamma/4$.
Choose $L_1$ sufficiently large that, when
$a>L_1R_+$, subtracting $\lambda_2R_+^{-4}$ from either upper cap
places it strictly below $b_{R_+}$. This is possible since
$|b_{R_+}|\le Ca^{-4}$ there. The preceding probability estimates
now apply with this fixed $L_1$.

Choose a safety constant $e_0>0$ so small that
$\eta=e_0R_+^{-4}$ is smaller than the barrier margin above $h_b$
throughout $R_+\le a\le2L_1R_+$ and than the shift gap. Finally choose
$\delta_0>0$ small compared with that margin, $\lambda_2$, and
$\lambda_1-\lambda_2$, so that
\begin{equation}\label{m:eq:cond-shift-order}
 16\delta_0+2e_0<\lambda_1-\lambda_2,\qquad
 \delta_0<\lambda_2,\qquad 2e_0<c'(2L_1)^{-4}.
\end{equation}
These choices precede the final smallness requirement on $r_*$.

Define
\[
 T_y(h)=\frac{f_y(h)}{4\pi},\qquad
 I_y(m)=\begin{cases}
 h_b(y),&m\le T_y(h_b(y)),\\
 T_y^{-1}(m),&T_y(h_b(y))<m<T_y(h_{\rm top}(y)),\\
 h_{\rm top}(y),&m\ge T_y(h_{\rm top}(y)).
 \end{cases}
\]
The positivity on the master kernel established with the barriers
implies
\begin{equation}\label{m:eq:cond-inverse-derivative}
 T_y'(h)\ge ca_y^{-2}
       \quad(h\ge h_b(y)-\delta_0a_y^{-4}).
\end{equation}
Indeed, differentiating the power $3/2$ leaves a positive square root
bounded below by $ca_y^{-2}$ on a kernel of integral $1+o(1)$.
Thus the inverse on the indicated interval exists and is continuous.
The clipped inverse satisfies
\begin{equation}\label{m:eq:cond-clipping}
 \begin{aligned}
 m\ge T_y(h_b)&\ \Longrightarrow\ T_y(I_y(m))\le m,\\
 m\le T_y(h_{\rm top})&\ \Longrightarrow\ T_y(I_y(m))\ge m.
 \end{aligned}
\end{equation}

\begin{lemma}[One-sided inverse comparisons]\label{m:lem:cond-inverse}
Fix $j<n$, and put $r=r_j$ and $R_+=2r$. For each deterministic $y$ with $r\le a_y\le2L_1R_+$, the failure
probability of each implication below is less than $r^{40}$:
\begin{equation}\label{m:eq:cond-inverse}
 \begin{aligned}
 \mu_j(y)\ge T_y(h_b(y))
   &\ \Longrightarrow\ H_j(y)\ge I_y(\mu_j(y))-\delta_0a_y^{-4},\\
 \mu_j(y)\le T_y(h_{\rm top}(y))
   &\ \Longrightarrow\ H_j(y)\le I_y(\mu_j(y))+\delta_0a_y^{-4}.
 \end{aligned}
\end{equation}
The constants and the smallness threshold for $r_*$ are independent
of the system and the number of scales.
\end{lemma}

\begin{proof}
The argument has three steps. Observation matching first relates
the conditional density to the particles retained by a fresh patch.
The patch comparison then determines the local offset except for a
small event, whose negative-field cost must also be controlled.
Finally, two separate tower identities transfer the offset comparison
back to the master field.

\paragraph{A count-good event and the fresh patch.}
Write $a=a_y$, $t=t_y$, $\tau=t/a$. Recall that
$Rf(y)=\int\mathcal K_z(y)f(z)\,dz$: the averaging and the kernel
gradient estimates below concern the center variable $z$.
There is a universal $K_0$ such that
\begin{equation}\label{m:eq:cond-observed-count}
 \Prob\{\#\{i:Y_{ji}\in B(y,3a)\}>K_0a^{-3}\}<r^{42}.
\end{equation}
Otherwise tilt this event. Lemma~\ref{lem:obs-tilt} gives
$\delta\le Cr^{-2.02}\log^5(e/r^{42})\le a^{-7+.02}$ for small
$r_*$. The last inequality is uniform: after multiplication by
$a^{6.98}$ its left side is at most
$C_{L_1}r^{4.96}\log^5(e/r^{42})$.
Every compatible raw configuration has the same excessive count in
$B(y,4a)$, because $\sqrt3\ell_j=o(a)$.
A fixed cell cover and Lemma~\ref{m:lem:screen-count} contradict this
for large enough $K_0$.

Call the complement of this excessive-count event count-good. On
such an observation path, match its entries to any compatible raw
configuration. The kernel Lipschitz bound gives
\begin{equation}\label{m:eq:cond-matching}
 \left|\sum_i\mathcal K_{x_i}(y)
             -\sum_i\mathcal K_{Y_{ji}}(y)\right|
       \le Ca^{-3}\ell_jt^{-4}.
\end{equation}
Only observations in $B(y,3a)$ can contribute to this difference.
The conditional ratios in Lemma~\ref{lem:obs-data} are supported on
compatible configurations. Averaging \eqref{m:eq:cond-matching}
therefore gives the same bound with $\mu_j(y)$ replacing the raw sum.

If an implication in \eqref{m:eq:cond-inverse} fails with probability
at least $r^{40}$, intersect its failure event with count-good data
and call the result $A$. Then $P_A:=\Prob(A)\ge r^{41}$ for small $r$.
Write $P$ for the original joint law of raw variables and all noisy
arrays, and $Q=P(\cdot\mid A)$. Its raw marginal is the law of
$\psi_A=\sqrt{p_A/P_A}\Psi$, with
$\delta\le a^{-7+.02}$.
Apply Lemma~\ref{m:lem:screen-patch} to this state, choosing
$b=a^{1+.2}$ and $\beta=b/a=a^{.2}$.
Once selected by shell averaging, the patch radius is deterministic.
Extend $Q$ by sampling the fresh cut labels conditionally on the raw
variables with the cut's squared partition probabilities.
The raw-label marginal is exactly the patch construction for
$\psi_A$. Write $X$ for its cut data, $\rho_X^c$ for its conditional
core density, $\rho$ for its local Thomas--Fermi minimizer, and
$\sigma$ for its retained fine density. These are conditional objects
under this tilted patch law. In contrast, $\mu_j$ remains the
conditional density defined under $P$.

\begin{center}
\small
\begin{tabular}{@{}p{.12\linewidth}p{.82\linewidth}@{}}
\toprule
Law & Construction and conditional quantities\\
\midrule
$P=\Prob$ & The raw state $\Psi$ with independent noisy arrays.
The fields $\mu_j,H_j$ retain their definitions using $\mathcal F_j$
under this law.\\[.6ex]
$Q$ & $P(\cdot\mid A)$, extended by the fresh cut labels.
The quantities $\rho_X^c,\rho,\sigma$ use its raw-label marginal
and the cut data $X$.\\
\bottomrule
\end{tabular}
\end{center}

The fine smearing moves retained centers by at most $\sqrt3b$.
Every particle relevant to the master kernel is in the retained
central region. Thus \eqref{m:eq:cond-matching} implies, $Q$-almost surely,
\begin{equation}\label{m:eq:cond-fine-matching}
 |R\sigma(y)-\mu_j(y)|\le Ca^{-3}(\ell_j+b)t^{-4}=o(a^{-6}).
\end{equation}
Let $G=\{D(\sigma-\rho)\le a^{-7+.01}\}$.
Lemma~\ref{m:lem:screen-density} gives
\begin{equation}\label{m:eq:cond-patch-budget}
 \E_Q\left[D(\sigma-\rho)+\int W_-\sigma\right]
       \le Ca^{-7+.02},\qquad
 \E_Q\int\sigma^{5/3}\le Ca^{-7}.
\end{equation}
In particular $Q(G^c)\le Ca^{.01}$. On $G$,
Lemma~\ref{lem:loc-coulomb} gives
\[
 |R\rho(y)-R\sigma(y)|\le Ct^{-5/2}a^{-3.5+.005}=o(a^{-6}).
\]
For clarity, after division by $a^{-6}$ these two errors are bounded
by
\begin{equation}\label{m:eq:cond-density-errors}
 C\{a^{.01}\tau^{-4}+a^{.2}\tau^{-4}
                         +a^{.005}\tau^{-5/2}\}=o(1).
\end{equation}
The lower bound $\tau\ge c_{L_1}a^w$, $w=10^{-5}$, makes every power
positive.

\paragraph{The local inverse comparison and rare negative fields.}
Put $h_X=W-V_c$. On $G$, in the high-density case one must have
\begin{equation}\label{m:eq:cond-good-high}
 h_X(y)\ge I_y(\mu_j(y))-\tfrac14\delta_0a^{-4}.
\end{equation}
Indeed, the opposite inequality and
Lemma~\ref{m:lem:barrier-oscillation} imply on the master kernel
$h_X(z)\le I_y(\mu_j(y))-\delta_0a^{-4}/8$.
To cover arbitrarily negative $h_X(y)$, note that
$h\mapsto h+C\tau h_-$ is increasing for small $\tau$; evaluate it
at the proposed upper threshold, which is bounded in $a^{-4}$ units.
The oscillation error there is $o(a^{-4})$.
The Thomas--Fermi equation, \eqref{m:eq:cond-inverse-derivative}, and
the first part of \eqref{m:eq:cond-clipping} then yield
$R\rho(y)\le\mu_j(y)-c\delta_0a^{-6}$, contradicting
\eqref{m:eq:cond-density-errors}.
In the low-density case the lower oscillation estimate and the
increasing function $h\mapsto h-C\tau h_-$ similarly give, on $G$,
\begin{equation}\label{m:eq:cond-good-low}
 h_X(y)\le I_y(\mu_j(y))+\tfrac14\delta_0a^{-4}.
\end{equation}
Here the second clipping inequality is used.

The deterministic upper bound $h_X(y)\le Ca^{-4}$ controls the
positive error on $G^c$ in the low-density case. The high-density case
requires a bound on negative fields there. By kernel positivity,
$\mu_j\ge T_y(h_b)\ge ca^{-6}$. Equation~\eqref{m:eq:cond-fine-matching}
and $\mathcal K_z(y)\le Ct^{-3}$ force, on every configuration,
\[
 m_\sigma:=\int_{B(y,2t)}\sigma\ge ca^{-6}t^3=ca^{-3}\tau^3.
\]
The upper oscillation bound implies on this ball
\[
 W_-\ge(1-C\tau)h_X(y)_--V_c-Ca^{-4}\tau^{1/2}.
\]
Integrate against $\sigma$, divide by $m_\sigma$, and average over
$G^c$. The patch bound controls the $W_-\sigma$ numerator by
$Ca^{-7+.02}$. For the nuclear term, weighted convexity and the
uniform integral of $|z-R_m|^{-5/2}$ on a radius-$2t$ ball give
\[
 \int_{B(y,2t)}V_c^{5/2}\le Ca^{-7}\tau^{1/2}.
\]
H\"older's inequality on probability times space therefore gives
\[
 \begin{split}
 \E_Q\1_{G^c}\int_{B(y,2t)}V_c\sigma
 &\le\left(\E_Q\int\sigma^{5/3}\right)^{3/5}
       \left(Q(G^c)\int_{B(y,2t)}V_c^{5/2}\right)^{2/5}\\
 &\le Ca^{-7+.004}\tau^{.2}.
 \end{split}
\]
We conclude that
\begin{equation}\label{m:eq:cond-negative-error}
 \E_Q\1_{G^c}h_X(y)_-
 \le Ca^{-4}\{a^{.02}\tau^{-3}+a^{.004}\tau^{-2.8}
                              +Q(G^c)\tau^{1/2}\}=o(a^{-4}).
\end{equation}
Since $|I_y|\le Ca^{-4}$, the bounds on $G$ now imply the
corresponding one-sided comparisons in expectation, with error
$\delta_0a^{-4}/4+o(a^{-4})$.

\paragraph{Transfer of averaged fields.}
The remaining step compares the expectation of this fresh patch
offset with that of the master offset under the event law $Q$.
Apply Lemma~\ref{lem:obs-transfer} with the nuclear splitting
$V=V_s+V_c$, master width $t=t_y$, and fine width $b=a^{1.2}$.
Lemma~\ref{m:lem:screen-density} gives the raw local density and Coulomb
error bounds; the deleted shell estimate is
\eqref{m:eq:screen-deleted-shell}. The pointwise density bound required
for the nuclear remainder follows from
$\rho=kW_+^{3/2}$ and \eqref{m:eq:screen-interior-upper}; its potential
on a radius-$t$ ball is $Ca^{-4}(t/a)^{1/2}$ by
\eqref{m:eq:local-power-potential}.
Thus the explicit transfer estimate is
\begin{equation}\label{m:eq:cond-transfer}
 |\E_Q(H_j-h_X)(y)|
 \le Ca^{-4}\{\tau^{1/5}+\beta^{1/5}+\beta^{1/2}
                +a^{.01}\tau^{-1/2}+\tau^{1/2}\}=o(a^{-4}).
\end{equation}
Here $\beta=a^{.2}$, $c_{L_1}a^w\le\tau\le c_1r_*^w$, and
$.01-w/2>0$. The two expectations in the common lemma use different
conditional laws; this verification identifies neither posterior pathwise.

On the high failure event, every path has
$H_j(y)<I_y(\mu_j(y))-\delta_0a^{-4}$.
Equations~\eqref{m:eq:cond-good-high} and~\eqref{m:eq:cond-negative-error},
together with \eqref{m:eq:cond-transfer}, instead give
\[
 \E_QH_j(y)\ge\E_QI_y(\mu_j(y))
                 -\tfrac14\delta_0a^{-4}-o(a^{-4}),
\]
a contradiction. The low failure event is contradicted in the same
way by \eqref{m:eq:cond-good-low}, the deterministic upper bound on
$h_X$, and \eqref{m:eq:cond-transfer}. All remainders are uniform:
$a\asymp_{L_1}r$, $\tau\ge c_{L_1}a^w$, and the smallest power
in \eqref{m:eq:cond-negative-error} is $.004-2.8w>0$.
The constants were fixed before letting $r_*\to0$.
\end{proof}

\subsection{Moments for the integrated exceptional cost}

The inverse estimates concern each deterministic spatial point.
Their use in the next section requires integrable costs for failed
comparisons, not a single event on which all points succeed.

\begin{lemma}[Conditional-density and reaction moments]\label{m:lem:cond-moments}
Fix $j<n$, put $r=r_j$ and $R_+=2r$, and let
$r\le a_y\le2L_1R_+$. Under the original joint law $P=\Prob$,
\begin{equation}\label{m:eq:cond-field-moments}
 \|\mu_j(y)\|_2\le Ct_y^{-3}a_y^{-3},\qquad
 \|f_y(H_j(y))\|_2
        \le Ca_y^{-6}(1+\tau_y^{-3/2}).
\end{equation}
Consequently, if $F_y$ is either failure event from
Lemma~\ref{m:lem:cond-inverse}, then
\begin{equation}\label{m:eq:cond-integrated-cost}
 \int_{\{r\le a_y\le2L_1R_+\}}
 \E\big[\1_{F_y}(\mu_j(y)+f_y(H_j(y)))\big]\dd y
       \le CM r^{20-3-3w}.
\end{equation}
The same bound holds with $F_y$ replaced by $\{y\in\mathcal E\}$,
and the expected integral over all of $\R^3$ of
$r^{-2}\sum_QB_Q\1_{Q'}$ is bounded by the right side as well.
\end{lemma}

\begin{proof}
Conditional Jensen, the master-kernel support, and
Lemma~\ref{m:lem:screen-count} give
\[
 \|\mu_j(y)\|_2\le Ct_y^{-3}
       \|\#\{i:x_i\in B(y,a_y)\}\|_2
       \le Ct_y^{-3}a_y^{-3}.
\]
The tail bound implies $\|(H_j(y)_+)^{3/2}\|_2\le Ca_y^{-6}$.
For the nuclear part of the reaction, weighted convexity and the
uniform integral of $|z-R_m|^{-3/2}$ on a radius-$2t_y$ ball give
\[
 \int\mathcal K_z(y)V_c(z)^{3/2}\dd z
       \le Ct_y^{-3}a_y^{-9/2}t_y^{3/2}
       =Ca_y^{-6}\tau_y^{-3/2}.
\]
This proves \eqref{m:eq:cond-field-moments}.
Cauchy--Schwarz against $\Prob(F_y)<r^{40}$ contributes $r^{20}$.
Use $a_y\asymp_{L_1}r$, $\tau_y\ge c_{L_1}r^w$, and
\eqref{m:eq:cond-band-volume} to obtain
\eqref{m:eq:cond-integrated-cost}. The arbitrarily high polynomial
bounds in Lemma~\ref{m:lem:cond-cap} give the other assertions, choosing
their exponent larger if necessary. All integrations are licensed
by the joint measurable versions already constructed.
\end{proof}
